% Generated by Markdown Studio v8 - ΣΗΜΜΥ solution template
% Template ID: shmmy-purple-cover-v1
% WARNING: review the exported equations against the approved source before publishing.
\documentclass[a4paper,11pt]{extarticle}
\usepackage{fontspec}
\usepackage{polyglossia}
\setmainlanguage{greek}
\setotherlanguage{english}
\IfFontExistsTF{DejaVu Serif}{\setmainfont{DejaVu Serif}}{\setmainfont{FreeSerif}}
\IfFontExistsTF{DejaVu Sans}{\setsansfont{DejaVu Sans}}{\setsansfont{FreeSans}}
\usepackage{amsmath,amssymb,graphicx,xcolor,geometry,enumitem,fancyhdr,hyperref}
\usepackage{longtable,booktabs,array}
\providecommand{\tightlist}{\setlength{\itemsep}{0pt}\setlength{\parskip}{0pt}}
\geometry{a4paper,margin=23mm}
\definecolor{accent}{HTML}{7446a7}
\hypersetup{colorlinks=true,linkcolor=accent,urlcolor=accent,pdftitle={Γραμμική Άλγεβρα}}
\setlength{\parindent}{0pt}
\setlength{\parskip}{5pt}
\setcounter{secnumdepth}{0}
\usepackage{titlesec}
\titleformat{\section}{\large\bfseries\color{accent}}{}{0pt}{}
\titleformat{\subsection}{\normalsize\bfseries}{}{0pt}{}
\titleformat{\subsubsection}{\normalsize\bfseries}{}{0pt}{}
\pagestyle{fancy}
\fancyhf{}
\fancyfoot[C]{\footnotesize\thepage}
\renewcommand{\headrulewidth}{0pt}
\begin{document}

\begin{titlepage}
\centering
\vspace*{2.4cm}
{\Large\bfseries\color{accent} ΟΜΑΔΑ ΕΠΙΛΥΣΗΣ ΘΕΜΑΤΩΝ ΣΗΜΜΥ\par}
\vspace{1.4cm}
{\Huge\bfseries Γραμμική Άλγεβρα\par}
\vspace{0.7cm}
{\Large Επαναληπτική Εξέταση -- 2025\par}
\vspace{2cm}
\begin{tabular}{@{}rl@{}}
\textbf{Λύτης:} & Name \\
\textbf{Ελεγκτής Α:} & Name1 \\
\textbf{Ελεγκτής Β:} & Name2 \\
\end{tabular}
\vfill
{\normalsize\itshape Λύσεις γραμμένες από φοιτητές για φοιτητές\par}
\vspace{1cm}
\end{titlepage}

\section{Άσκηση 1 α)}\label{ux3acux3c3ux3baux3b7ux3c3ux3b7-1-ux3b1}

\textbf{Μεθοδολογία}

\begin{enumerate}
\def\labelenumi{\arabic{enumi}.}
\tightlist
\item
  Υπολογίζουμε την ισοδύναμη αντίσταση.
\item
  Εφαρμόζουμε τον νόμο του Ohm.
\end{enumerate}

\textbf{Λύση}

\[
R_{\mathrm{eq}} = R_1+R_2 = 20\,\Omega
\]

\[
I=\frac{100}{20}=5\,\mathrm{A}
\]

Επομένως, το ρεύμα είναι \(I=5\,\mathrm{A}\).

\section{Άσκηση 1 β)}\label{ux3acux3c3ux3baux3b7ux3c3ux3b7-1-ux3b2}

\textbf{Μεθοδολογία}

Χρησιμοποιούμε το αποτέλεσμα του α) ερωτήματος.

\textbf{Λύση}

Η ισχύς προκύπτει από την σχέση \(P=I^2R\).

\end{document}
